\documentclass{article}

\usepackage{listings}
\usepackage{color}
\usepackage{graphicx}
\usepackage{float}
\usepackage{amsmath}
\usepackage{subfig}
\usepackage{cite}
\usepackage{url}
\usepackage{amsmath}

\newcounter{qcounter}

\begin{document}

\title{Image Analysis - TP5 - Lucas Kanade Tracker}

\author{Jander Nascimento, 
\and Raquel Oliveira}

\maketitle

\section{Initial concept}


\section{Interpolation}

	Bilinear Interpolation is a method to generate a new image (given a initial image {\it I}) where the value of the pixels is a sum of the intensities of its four neighbors, weighted by the distance to the neighbor's center.

	For each pixel {\it(i,j)} of the original image {\it I}, the new pixel in the interpolated image {\it $I_w$} is calculated in the following manner:

	\begin{equation}
		I_w(i+\Delta_i^e , j+\Delta_j^e) += (1-\Delta_i^f)(1-\Delta_j^f)I(i,j)
	\end{equation}
	\begin{equation}
		I_w(i+\Delta_i^e+1 , j+\Delta_j^e) += \Delta_i^f(1-\Delta_j^f)I(i,j)
	\end{equation}
	\begin{equation}
  		I_w(i+\Delta_i^e , j+\Delta_j^e+1) += (1-\Delta_i^f)\Delta_j^fI(i,j)
	\end{equation}
	\begin{equation}
		I_w(i+\Delta_i^e+1 , j+\Delta_j^e+1) += \Delta_i^f \Delta_j^fI(i,j)
	\end{equation}

	where: 
		\begin{itemize}
			\item {\it $\Delta_i$} is the translation along {\it i}
			\item {\it $\Delta_j$} is the translation along {\it j} 
			\item {\it $\Delta^e$} is the integer part of the translation (ie: $\Delta=-5.3$, then ${\it \Delta^e}=-5$
			\item {\it $\Delta^f$} is the decimal part of the translation (ie: $\Delta=-5.3$, then ${\it \Delta^f}=0.3$
		\end{itemize}

	In the Figure \ref{fig:interp} we can see the result of an interpolation. It is possible to observe that the interpolated image, compared with the original image, is a little deformed, due to the fact that the pixels in the new image are interpolated with their neighbors.

	\begin{figure}[H]
		  \centering
		  \subfloat[Original image]{\includegraphics[width=0.4\textwidth]{img/taz.png}}
		  \hspace{0.1cm}
		  \subfloat[Interpolated image ($\Delta_i=12.8   \Delta_j=-6.3$)]{\includegraphics[width=0.4\textwidth]{img/taz0interp.png}}
		  \caption{Interpolation of an image}
		  \label{fig:interp}
	\end{figure}	


\section{Anaysis of a static image}


\section{Tranking an object}


\section{How to run?}

	Steps to compile the application:
	
	\begin{itemize}
		\item svn checkout https://jfimageanalysis.googlecode.com/svn/trunk/TP4/ \#download source code
		\item make \#compiles the code
	\end{itemize}

	As an input image only {\bf ppm plaintext/ansii} files are accepted (P3). 

	Examples of usage:

	\begin{itemize}
		\item Create 3 clusters:
		\subitem ./showregion -i image\_in.ppm -g 3 $>$ image\_out.ppm
	\end{itemize}

	You can always type {\it ./tp5 --help} to check for more options. 

\end{document}


